\section{Discussion}\label{sec:discussion}

\subsection{What does structural similarity tell us?}
PN-GDDA describes similarity of local directed bipartite Petri-net wiring.
The GIM value 0.9976 is therefore meaningful but deliberately incomplete as a
statement about executable behavior. It remains constant when labels change
because the underlying structure does not. A label-blind score cannot determine
labeled trace equality or simulation in general (Proposition
\ref{prop:label-blind}); this is a boundary of the question it answers, not a
defect in the comparator. Structural and behavioral comparisons are
complementary, rather than competitors for one universal notion of similarity.

\subsection{What does behavioral comparison add?}
The executable comparison determines whether observable steps can be matched
recursively from the declared initial states. Unlike a trace list, it retains
the location at which choices become committed (Example~\ref{ex:branching}).
This distinction is relevant when transferring qualitative model behavior
requires preserving alternative continuations rather than merely permitting
the same sequences. Related formal work on asynchronous Boolean networks also
emphasizes reachable branching in phenotypic analysis.\cite{Yizengaw2026}

In GIM, the model-level answer is correspondence at A/B and its loss at C.
In RCD, only the plant model's declared branching behavior is reproducible by
the animal model, not conversely. Simulation is stronger than saying that
each plant trace occurs in the animal model: every related successor pair
must continue to satisfy the matching condition. Neither result transfers to
other interfaces or organisms without re-analysis.

\subsection{What does interface refinement mean biologically?}
An observational interface states which differences a biological readout is
allowed to retain, merge, or hide. At A/B, the GIM terminal responses are
collapsed into loss of proliferative capacity; at C, their encoded mechanisms
are distinguishable. Their incompatibility at C is consequently consistent
with the model construction, not independent evidence that the method has
discovered a new organism-level difference. The contribution is to determine
the survival or failure of recursive correspondence for the full executable
models under each explicit interface.

The result localizes a boundary within three declared readouts; it does not
optimize over interfaces or identify a unique minimal observational scale.
Because the interfaces were informed by known mechanisms, the C result is not
a blinded prediction. A prospective test would need independently selected
readouts, matched biological conditions, and outcomes not used in model or
interface construction. The present assay interpretation remains a hypothesis.
The distinction between a combined differentiation/endoreduplication transition
and two independently selectable outcomes is also essential: only the combined
transition is encoded here.

\subsection{Execution assumptions and evidence limits}
The comparison depends on transitions, initial states, labels, and update
semantics, not on graph structure alone. Asynchronous versus synchronous
execution changes public-model reachability and two HPN classes. General
asynchronous, most-permissive, and stochastic semantics were not evaluated.
\cite{Pauleve2020} Weak bisimulation here is divergence-insensitive and concerns
untimed interleaving behavior, not rates, probabilities, or true concurrency.
Controlled silent refinement preserves the relation only under the stated
construction (Proposition~\ref{prop:silent-refinement}), not under arbitrary
insertions of conflicts or observable branches.

The GIM pair is manually curated and compact, with no kinetics, concentrations,
tissue context, or independent experimental calibration. Literature justifies
the chosen biological distinctions but does not certify the entire executable
model. HPN provides independent measurements yet has only nine pair--condition
observations, three cell lines, and inconclusive associations. Its formal
directions are nominal; assigning them an ordinal score would lose information.
CASPOTS RMSE is minimum-distance compatibility over admissible trajectories,
not a unique time-series prediction. These limitations cannot be removed by
adding more software controls.

\subsection{Computational scope}
Explicit reachability remains vulnerable to state-space explosion, and a
candidate relation stores the product of the two state sets. The scaling study
in Supplementary Validation is illustrative, not evidence of genome-scale
performance. Applying native PN-GDDA to logical rules would additionally require
a non-unique Petri-net translation. Larger independently curated model pairs,
prospectively chosen interfaces, and empirical readouts are needed before
claiming wider biological utility. The present method is most defensible as a
reproducible, question-specific comparison of explicit qualitative models.

\section{Conclusion}\label{sec:conclusion}

Two curated DNA-damage response models with constant PN-GDDA similarity 0.9976
retain recursive behavioral correspondence at coarse and pathway resolution,
but become non-comparable when their encoded terminal mechanisms are separately
observed. For these models, correspondence depends on observational resolution;
the boundary lies between Interfaces B and C within the declared family.

The computational contribution is to make that conditional boundary explicit
and reproducible, while preserving the difference between similar wiring,
matching traces, and matching branching. Proofs and independent checks support
the formal decisions without substituting for biological evidence. The result
does not establish organism-level equivalence, mechanistic identity, or
prognosis. Its empirical relevance remains to be tested beyond these curated
models; the secondary HPN analysis is exploratory and inconclusive.
